\documentclass[10pt,a4paper]{amsart}
\usepackage{amsmath,amssymb,amsthm,graphicx}
% For commutative diagrams we use Xy-pic
\usepackage[all]{xy}

% We use 2cell for 2-commutative diagrams.
\xyoption{2cell}
% Theorem environments.
%
\theoremstyle{plain}
\newtheorem{theorem}[subsection]{Theorem}
\newtheorem{proposition}[subsection]{Proposition}
\newtheorem{lemma}[subsection]{Lemma}

\theoremstyle{definition}
\newtheorem{definition}[subsection]{Definition}
\newtheorem{example}[subsection]{Example}
\newtheorem{exercise}[subsection]{Exercise}
\newtheorem{situation}[subsection]{Situation}

\theoremstyle{remark}
\newtheorem{remark}[subsection]{Remark}
\newtheorem{remarks}[subsection]{Remarks}

\numberwithin{equation}{subsection}

% Macros
%
\def\lim{\mathop{\mathrm{lim}}\nolimits}
\def\colim{\mathop{\mathrm{colim}}\nolimits}
\def\Spec{\mathop{\mathrm{Spec}}}
\def\Hom{\mathop{\mathrm{Hom}}\nolimits}
\def\Ext{\mathop{\mathrm{Ext}}\nolimits}
\def\SheafHom{\mathop{\mathcal{H}\!\mathit{om}}\nolimits}
\def\SheafExt{\mathop{\mathcal{E}\!\mathit{xt}}\nolimits}
\def\Sch{\mathit{Sch}}
\def\Mor{\mathop{\mathrm{Mor}}\nolimits}
\def\Ob{\mathop{\mathrm{Ob}}\nolimits}
\def\Sh{\mathop{\mathit{Sh}}\nolimits}
\def\NL{\mathop{N\!L}\nolimits}
\def\CH{\mathop{\mathrm{CH}}\nolimits}
\def\proetale{{pro\text{-}\acute{e}tale}}
\def\etale{{\acute{e}tale}}
\def\QCoh{\mathit{QCoh}}
\def\Ker{\mathop{\mathrm{Ker}}}
\def\Im{\mathop{\mathrm{Im}}}
\def\Coker{\mathop{\mathrm{Coker}}}
\def\Coim{\mathop{\mathrm{Coim}}}

% Boxtimes
%
\DeclareMathSymbol{\boxtimes}{\mathbin}{AMSa}{"02}

%
% Macros for moduli stacks/spaces
%
\def\QCohstack{\mathcal{QC}\!\mathit{oh}}
\def\Cohstack{\mathcal{C}\!\mathit{oh}}
\def\Spacesstack{\mathcal{S}\!\mathit{paces}}
\def\Quotfunctor{\mathrm{Quot}}
\def\Hilbfunctor{\mathrm{Hilb}}
\def\Curvesstack{\mathcal{C}\!\mathit{urves}}
\def\Polarizedstack{\mathcal{P}\!\mathit{olarized}}
\def\Complexesstack{\mathcal{C}\!\mathit{omplexes}}
% \Pic is the operator that assigns to X its picard group, usage \Pic(X)
% \Picardstack_{X/B} denotes the Picard stack of X over B
% \Picardfunctor_{X/B} denotes the Picard functor of X over B
\def\Pic{\mathop{\mathrm{Pic}}\nolimits}
\def\Picardstack{\mathcal{P}\!\mathit{ic}}
\def\Picardfunctor{\mathrm{Pic}}
\def\Deformationcategory{\mathcal{D}\!\mathit{ef}}
\usepackage{cleveref}
\crefname{originallemma}{original lemma}{original lemmas}
\Crefname{originallemma}{Original lemma}{Original lemmas}
\makeatletter
\newlabel{flat-lemma-refine-check-h-stack}{{0EX9}{source}{}{sourcepointer-1}{}}
\newlabel{flat-lemma-refine-check-h-stack@cref}{{[originallemma][0][]0EX9}{[1][0][]source}{}{}{}}
\makeatother
\crefname{originalreference}{original reference}{original references}
\Crefname{originalreference}{Original reference}{Original references}
\makeatletter
\newlabel{lemma-glue-etale-sheaf-fppf-cover}{{0GF2}{source}{}{sourcepointer-2}{}}
\newlabel{lemma-glue-etale-sheaf-fppf-cover@cref}{{[originalreference][0][]0GF2}{[1][0][]source}{}{}{}}
\makeatother
\crefname{originalreference}{original reference}{original references}
\Crefname{originalreference}{Original reference}{Original references}
\makeatletter
\newlabel{lemma-glue-etale-sheaf-modification}{{0GF4}{source}{}{sourcepointer-3}{}}
\newlabel{lemma-glue-etale-sheaf-modification@cref}{{[originalreference][0][]0GF4}{[1][0][]source}{}{}{}}
\makeatother
\crefname{originalreference}{original reference}{original references}
\Crefname{originalreference}{Original reference}{Original references}
\makeatletter
\newlabel{lemma-glue-etale-sheaf-proper-surjective}{{0GF0}{source}{}{sourcepointer-4}{}}
\newlabel{lemma-glue-etale-sheaf-proper-surjective@cref}{{[originalreference][0][]0GF0}{[1][0][]source}{}{}{}}
\makeatother
\crefname{originallemma}{original lemma}{original lemmas}
\Crefname{originallemma}{Original lemma}{Original lemmas}
\makeatletter
\newlabel{sites-lemma-glue-sheaves}{{04TR}{source}{}{sourcepointer-5}{}}
\newlabel{sites-lemma-glue-sheaves@cref}{{[originallemma][0][]04TR}{[1][0][]source}{}{}{}}
\makeatother
\title{The stack of etale sheaves on schemes over a base for the h-topology}
\author{The Stacks Project Authors}
\date{stacks-project:etale-cohomology.tex:local-snapshot-20261011; GNU FDL 1.2 or later}
\begin{document}
\maketitle
\noindent\textit{Editorial note.} All delivered mathematical text is the author's own, copied from the pinned Stacks Project chapter file etale-cohomology.tex; no statement or proof step was rewritten, reordered or removed.\par\medskip
\noindent\textit{Editorial note.} The counted proof is the author's complete argument: the verification of the three conditions of the h-stack criterion, namely Zariski glueing of sheaves, descent along an fppf cover or a proper surjective cover, and the modification condition.\par\medskip
\noindent\textit{Editorial note.} The label lemma-h-descent-etale-sheaves names the counted statement, and the proof environment immediately following it in the source is delivered verbatim as its proof.\par\medskip
\noindent\textit{Editorial note.} The counted statement is carried by the author's own lemma environment in the chapter source.\par\medskip
\noindent\textit{Editorial note.} Pre-open gate: python operations/agent-authoring/tools/precard\_gate.py --source etale-cohomology.tex --label lemma-h-descent-etale-sheaves --start 23256 --end 23285 --local-id 3574 returned RESULT=OK (L3574 free, occupancy FREE, public interval clear). Candidate taken from operations/pool-cards/stacks-etale-cohomology-halfB.json.\par\medskip
\noindent\textit{Editorial note.} The statement and proof are copied from the pinned chapter file; the cited inputs are kept as external pointers with their canonical Stacks tags.\par\medskip
\noindent\textit{Allowed standard prerequisites (editorial).}\begin{itemize}
\item Criterion for a category fibred in groupoids to be an h-stack. More on Flatness, the lemma listing the three conditions (Zariski glueing, descent along surjective flat proper finitely presented morphisms, and the modification condition) that imply that a category fibred in groupoids over the h-site is a stack. Supplies the three conditions checked in the proof.
\item Glueing of sheaves and descent for an fppf and for a proper surjective cover. Sites, the lemma on glueing of sheaves; Same chapter, the lemma giving descent of etale sheaves along an fppf cover and the lemma giving descent along a proper surjective cover. Supply condition (1) and condition (2) respectively.
\item Descent of etale sheaves along a modification. Same chapter, the lemma establishing descent for the modification condition. Supplies condition (3).
\end{itemize}\medskip
\ifdefined\hypertarget\hypertarget{sourcepointer-1}{}\fi
\noindent\textit{Source reference.} Original lemma 0EX9: sources/stacks-chow/tags.txt, source lines 18649--18649. Prerequisite pointer invoked at proof line(s) 23268; it names a standard fact of a sibling chapter that the counted argument uses, and every step of the counted proof built on it is retained verbatim in the delivered body.\par\medskip
\ifdefined\hypertarget\hypertarget{sourcepointer-2}{}\fi
\noindent\textit{Source reference.} Original reference 0GF2: sources/stacks-etale-cohomology/pinned-original-source.tex, source lines 23100--23100. Prerequisite pointer invoked at proof line(s) 23279; it names a standard fact of the same chapter that the counted argument uses, and every step of the counted proof built on it is retained verbatim in the delivered body.\par\medskip
\ifdefined\hypertarget\hypertarget{sourcepointer-3}{}\fi
\noindent\textit{Source reference.} Original reference 0GF4: sources/stacks-etale-cohomology/pinned-original-source.tex, source lines 23167--23167. Prerequisite pointer invoked at proof line(s) 23284; it names a standard fact of the same chapter that the counted argument uses, and every step of the counted proof built on it is retained verbatim in the delivered body.\par\medskip
\ifdefined\hypertarget\hypertarget{sourcepointer-4}{}\fi
\noindent\textit{Source reference.} Original reference 0GF0: sources/stacks-etale-cohomology/pinned-original-source.tex, source lines 23047--23047. Prerequisite pointer invoked at proof line(s) 23280; it names a standard fact of the same chapter that the counted argument uses, and every step of the counted proof built on it is retained verbatim in the delivered body.\par\medskip
\ifdefined\hypertarget\hypertarget{sourcepointer-5}{}\fi
\noindent\textit{Source reference.} Original lemma 04TR: sources/stacks-chow/tags.txt, source lines 6264--6264. Prerequisite pointer invoked at proof line(s) 23273; it names a standard fact of a sibling chapter that the counted argument uses, and every step of the counted proof built on it is retained verbatim in the delivered body.\par\medskip
\begin{lemma}
\label{lemma-h-descent-etale-sheaves}
Let $S$ be a scheme. Then the category fibred in groupoids
$$
p : \mathcal{S} \longrightarrow (\Sch/S)_h
$$
whose fibre category over $U$ is the category $\Sh(U_\etale)$
of sheaves on the small \'etale site of $U$ is a stack in groupoids.
\end{lemma}

\begin{proof}
To prove the lemma we will check conditions (1), (2), and (3) of
More on Flatness, Lemma \ref{flat-lemma-refine-check-h-stack}.

\medskip\noindent
Condition (1) holds because we have glueing for sheaves (and
Zariski coverings are \'etale coverings). See
Sites, Lemma \ref{sites-lemma-glue-sheaves}.

\medskip\noindent
To see condition (2), suppose that $f : X \to Y$ is a surjective,
flat, proper morphism of finite presentation over $S$ with $Y$ affine.
Then we have descent for $\{X \to Y\}$ by either
Lemma \ref{lemma-glue-etale-sheaf-fppf-cover} or
Lemma \ref{lemma-glue-etale-sheaf-proper-surjective}.

\medskip\noindent
Condition (3) follows immediately from the more general
Lemma \ref{lemma-glue-etale-sheaf-modification}.
\end{proof}

\end{document}
